% !TeX root = main.tex
\section{Related Work}
Duty-cycle and split-cycle concepts have been explored as a way to reshape wing velocity profiles within a cycle for flapping-wing vehicles.
These approaches are attractive because they can be implemented as time reparameterizations of existing periodic trajectories without modifying mechanical stroke amplitude.
Quasi-steady models remain a common tool for interpreting trends and designing controllers due to their computational efficiency and their ability to produce cycle-averaged force and power estimates.

This paper builds on strip-theory quasi-steady modeling approaches for flapping wings and focuses on a control-oriented evaluation of duty-cycle modulation near trim-like conditions.
We adopt a phase-warped SCCPFM waveform inspired by prior experimental control-system work and use a DeLaurier-style quasi-steady strip-theory model for aerodynamic loading.
While the absolute accuracy of quasi-steady models can be limited in separated/unsteady regimes, they provide a tractable platform to quantify duty-cycle tradeoffs and to identify promising operating points for subsequent hardware evaluation.
